\section*{Capítulo \ref{ch:b3:colloids} --- Interfases, tensioactivos y coloides}

\begin{solution}{exo:b3:colloids:1}
$r = \qty{0.50}{\micro m}$: $\Delta p = 2 \times 0.0720/\num{5.0e-7} = \qty{2.9e5}{Pa}$; en el interior, $1.0 + 2.9 = \qty{3.9}{bar}$.
\end{solution}

\begin{solution}{exo:b3:colloids:2}
$\ln(p/p^*) = 2 \times 0.0720 \times \num{1.807e-5}/(\num{1.0e-8} \times 8.314 \times 298.15) = 0.105$; $p/p^* = 1.11$.
\end{solution}

\begin{solution}{exo:b3:colloids:3}
$\cos\theta = (40 - 20)/72 = 0.278$, $\theta = \ang{74}$: por debajo de $90^\circ$, el líquido \omterm{def:b3:colloids:wetting}{moja} el sólido
parcialmente.
\end{solution}

\begin{solution}{exo:b3:colloids:4}
NaCl: $I = \qty{10}{mol.m^{-3}}$, $\kappa^{-1} = 0.96\sqrt{100/10} = \qty{3.0}{nm}$. \ce{MgSO4}: $I = \qty{40}{mol.m^{-3}}$, \qty{1.5}{nm}.
\end{solution}

\begin{solution}{exo:b3:colloids:5}
$\Gamma = \num{12.6e-3}/(2 \times 8.314 \times 298.15) = \qty{2.54e-6}{mol.m^{-2}}$; área $1/(\Gamma N_A) =
\qty{0.65}{nm^2}$ por molécula.
\end{solution}

\begin{solution}{exo:b3:colloids:6}
Por debajo de la CMC, la pendiente es de $-12.9$ \unit{mN/m} por unidad de $\log c$ (de $-5.0$ a $-4.0$); la
meseta está en 33.4. La recta la alcanza en $\log c = -4.0 + (39.1 - 33.4)/12.9 = -3.56$: CMC
$\approx \qty{2.8e-4}{mol/L}$.
\end{solution}

\begin{solution}{exo:b3:colloids:7}
SDS: $P = 0.350/(0.60 \times 1.67) = 0.35$, en el límite entre esferas y cilindros
cortos: \omterm{def:b3:colloids:micelle}{micelas} esféricas cerca de la CMC. Lípido: $P = 0.700/(0.65 \times 1.67) = 0.64$:
bicapas, y por tanto vesículas.
\end{solution}

\begin{solution}{exo:b3:colloids:8}
\ce{CaCl2}: $60/64 = \qty{0.94}{mM}$ de \ce{Ca^{2+}}; \ce{AlCl3}: $60/729 = \qty{0.082}{mM}$ de \ce{Al^{3+}}. El \ce{Na3PO4}
coagula a través de su \ce{Na+} (a unos \qty{20}{mM} de sal); el fosfato, un
coión, puede incluso adsorberse y aumentar la carga negativa.
\end{solution}

\begin{solution}{exo:b3:colloids:9}
Parte hidrófila: $6 \times 44.0 + 17.0 = \qty{281.0}{g/mol}$ de \qty{450.0}{g/mol} (masas atómicas del libro):
HLB $= 20 \times 281.0/450.0 = 12.5$, un emulsionante de aceite en agua.
\end{solution}

\begin{solution}{exo:b3:colloids:10}
El aceite de las gotitas pequeñas es más soluble en la fase continua (por el
factor $\exp(2\gamma V_{\mathrm m}/rRT)$, cuyo exponente es diez veces mayor para \qty{50}{nm} que para \qty{500}{nm}): se
difunde hacia las gotitas grandes, que crecen mientras las pequeñas desaparecen. Esta
\omterm{def:b3:colloids:ripening}{maduración de Ostwald} engruesa una \omterm{def:b3:colloids:colloid}{emulsión} aunque nunca se fusionen dos gotitas.
\end{solution}

\begin{solution}{exo:b3:colloids:11}
El menisco es un casquete esférico de radio $r/\cos\theta$: el líquido justo por debajo de él está a una
presión menor en $2\gamma\cos\theta/r$, y la columna sube hasta que $\rho gh$ lo compensa. Para el agua,
$h = 2 \times 0.072/(997 \times 9.81 \times \num{1.0e-4}) = \qty{0.15}{m}$.
\end{solution}

\begin{solution}{exo:b3:colloids:12}
Con \qty{100}{mM}, la \omterm{def:b3:electrode-kinetics:double-layer}{longitud de Debye} es inferior a \qty{1}{nm} y la repulsión se extingue antes que la
atracción: no queda ninguna barrera (más allá de la curva de \qty{50}{mM}). Un polímero libre no
actúa en la superficie; una capa de polímero injertado repele al contacto sea cual sea la
sal, por la entropía que se pierde cuando las cadenas se solapan.
\end{solution}

\begin{solution}{pb:b3:colloids:1}
\textbf{1.} Las sustituciones en la red cristalina (\ce{Al^{3+}} por \ce{Si^{4+}}, \ce{Mg^{2+}} por \ce{Al^{3+}}) dejan
una carga negativa permanente, compensada por cationes intercambiables.
\textbf{2.} $I = \frac12(1.0 + 1.0) = \qty{1.0}{mM}$.
\textbf{3.} $\kappa^{-1} = \qty{9.6}{nm}$.
\textbf{4.} $a\kappa = 10$: la doble capa es delgada frente a la partícula, la
condición de la aproximación de Derjaguin.
\textbf{5.} $v = 2 \times 1650 \times 9.81 \times (\num{1.0e-7})^2/(9 \times \num{0.89e-3}) = \qty{4.0e-8}{m/s}$, unos \qty{3.5}{mm} al día.
\textbf{6.} Sus \omterm{def:b3:electrode-kinetics:double-layer}{capas difusas} se repelen antes de que se imponga la atracción de van der Waals.
\textbf{7.} $V = 2\pi\varepsilon_r\varepsilon_0a\psi^2\ln(1 + \eu^{-\kappa h}) - Aa/12h$.
\textbf{8.} Unos $39\,kT$.
\textbf{9.} Una fracción del orden de $\eu^{-39} \approx 10^{-17}$ de las colisiones: prácticamente nunca.
\textbf{10.} Desaparece: toda colisión conduce al contacto.
\textbf{11.} Una \omterm{def:b3:electrode-kinetics:double-layer}{longitud de Debye} más corta confina la repulsión a distancias pequeñas, donde la
atracción, que crece como $1/h$, domina.
\textbf{12.} La CCC varía como $z^{-6}$ de la carga del contraión.
\textbf{13.} $50/64 = \qty{0.78}{mM}$.
\textbf{14.} $50/729 = \qty{0.069}{mM}$.
\textbf{15.} Los cationes son los contraiones de las partículas negativas: se
acumulan en la doble capa y la apantallan.
\textbf{16.} El sulfato es un coión, repelido por las partículas.
\textbf{17.} \qty{0.069}{mol} de \ce{Al^{3+}} por metro cúbico.
\textbf{18.} \qty{0.034}{mol} de $\ce{Al2(SO4)3}$.
\textbf{19.} $0.0343 \times 342.3 = \qty{11.7}{g}$.
\textbf{20.} \qty{11.7}{kg} por hora.
\textbf{21.} $\ce{Al^3+ + 3 HCO3- -> Al(OH)3 + 3 CO2}$.
\textbf{22.} $3 \times 0.069 = \qty{0.21}{mM}$ de los \qty{1.0}{mM}: una quinta parte. El
hidrogenocarbonato restante y el \ce{CO2} disuelto tamponan el agua: el pH solo baja
moderadamente.
\textbf{23.} Las especies de aluminio muy cargadas se adsorben e invierten la carga de las
partículas, que vuelven a repelerse.
\textbf{24.} \textbf{Unos \qty{12}{g} de sulfato de aluminio por metro cúbico} (\qty{11.7}{g}), antes del suplemento
necesario en la práctica para la hidrólisis.
\end{solution}
